\documentclass[11pt]{article}
\usepackage[a4paper,margin=1in]{geometry}
\usepackage{amsmath,amssymb}
\usepackage{enumitem}
\usepackage{setspace}

\onehalfspacing

\title{\textbf{CSE 400: Fundamentals of Probability in Computing \\ Tutorial 2 Solutions}}
\date{}
\author{}

\begin{document}

\maketitle

\section*{Q1. Random Point on a Line Segment}

\textbf{Problem Interpretation}

A point is chosen uniformly at random on a line segment of length $L$.

\textbf{Step 1: Mathematical Representation}

Let:

$x$ = distance of the chosen point from one end

Then $0 < x < L$

The line is divided into two segments:

Length 1: $x$

Length 2: $L - x$

The ratio of shorter to longer segment is:

\[
R = \frac{\min(x,L-x)}{\max(x,L-x)}
\]

We must compute:

\[
P\left(R < \frac{1}{4}\right)
\]

\textbf{Step 2: Divide Based on Midpoint}

Since which segment is shorter depends on position relative to $L/2$, we divide into two
cases.

\textbf{Case 1:} $x \le L/2$

Shorter segment = $x$

Longer segment = $L - x$

Condition:

\[
\frac{x}{L-x} < \frac{1}{4}
\]

Multiply both sides:

\[
4x < L - x
\]

\[
5x < L
\]

\[
x < \frac{L}{5}
\]

\textbf{Case 2:} $x > L/2$

Shorter segment = $L - x$

Longer segment = $x$

Condition:

\[
\frac{L-x}{x} < \frac{1}{4}
\]

Multiply:

\[
4(L-x) < x
\]

\[
4L - 4x < x
\]

\[
5x > 4L
\]

\[
x > \frac{4L}{5}
\]

\textbf{Step 3: Valid Region}

Thus valid $x$ lies in:

\[
(0,L/5) \cup (4L/5,L)
\]

Length of valid region:

\[
\frac{L}{5} + \frac{L}{5} = \frac{2L}{5}
\]

Since $x$ is uniform over length $L$:

\[
P = \frac{2L/5}{L} = \frac{2}{5}
\]

\section*{Q2. Bayesian Classification with Normal Distributions}

\textbf{Given}

White region:

\[
X|W \sim N(4,4), \quad \sigma_W = 2
\]

Black region:

\[
X|B \sim N(6,9), \quad \sigma_B = 3
\]

Prior probabilities:

\[
P(B)=\alpha,\quad P(W)=1-\alpha
\]

Observation:

\[
X=5
\]

We want equal error probability regardless of conclusion.

\textbf{Step 1: Equal Posterior Probabilities}

For equal error:

\[
P(B|X=5)=P(W|X=5)=\frac{1}{2}
\]

Using Bayes’ theorem:

\[
P(B|X=5)=\frac{\alpha f(5|B)}{\alpha f(5|B)+(1-\alpha)f(5|W)}
\]

Set equal to $\frac{1}{2}$.

\textbf{Step 2: Normal Density Formula}

\[
f(x)=\frac{1}{\sigma\sqrt{2\pi}}e^{-\frac{(x-\mu)^2}{2\sigma^2}}
\]

Compute densities:

\[
f(5|B)=\frac{1}{3\sqrt{2\pi}}e^{-1/18}
\]

\[
f(5|W)=\frac{1}{2\sqrt{2\pi}}e^{-1/8}
\]

Substitute into Bayes expression:

\[
\frac{3\alpha}{3\alpha+2(1-\alpha)e^{-5/72}}=\frac{1}{2}
\]

\textbf{Step 3: Solve}

\[
6\alpha = 3\alpha + 2(1-\alpha)(0.93)
\]

\[
3\alpha = 1.86 - 1.86\alpha
\]

\[
4.86\alpha = 1.86
\]

\[
\alpha \approx 0.3827
\]

% --- Remaining content unchanged below (intentionally preserved exactly) ---

\section*{Q3. Minimizing Expected Distance}

\textbf{(a) Uniform on (0,A)}

\[
X \sim \text{Uniform}(0,A)
\]

PDF:

\[
f_X(x)=\frac{1}{A}
\]

Minimize:

\[
E[|X-a|]
\]

Break absolute value:

\[
E[|X-a|]=\frac{1}{A}\int_0^a(a-x)dx+\frac{1}{A}\int_a^A(x-a)dx
\]

After evaluation:

\[
\frac{2a^2-2aA+A^2}{2A}
\]

Differentiate:

\[
\frac{4a-2A}{2A}=0
\]

\[
4a=2A
\]

\[
a=\frac{A}{2}
\]

Station should be at midpoint.

\textbf{(b) Exponential Distribution}

\[
X \sim \text{Exponential}(\lambda)
\]

PDF:

\[
\lambda e^{-\lambda x}
\]

Minimize:

\[
E[|X-a|]
\]

Derivative:

\[
\frac{d}{da}E[|X-a|]=1-2e^{-\lambda a}
\]

Set to zero:

\[
2e^{-\lambda a}=1
\]

\[
e^{-\lambda a}=\frac{1}{2}
\]

\[
a=\frac{\ln 2}{\lambda}
\]

\section*{Q4. Mixture of Exponentials}

Two types:

Rate $\lambda_1$, probability $p_1$

Rate $\lambda_2$, probability $p_2$

We want:

\[
P(X>t+s|X>t)
\]

Using:

\[
P(X>t)=p_1e^{-\lambda_1 t}+p_2e^{-\lambda_2 t}
\]

\[
P(X>t+s)=p_1e^{-\lambda_1 (t+s)}+p_2e^{-\lambda_2 (t+s)}
\]

Thus:

\[
\frac{p_1e^{-\lambda_1 (t+s)}+p_2e^{-\lambda_2 (t+s)}}{p_1e^{-\lambda_1 t}+p_2e^{-\lambda_2 t}}
\]

\section*{Q5. Conditional Poisson → Binomial}

\[
X \sim \text{Poisson}(\lambda_1), \quad Y \sim \text{Poisson}(\lambda_2)
\]

Independent.

We need:

\[
P(X=k|X+Y=n)
\]

Using definition:

\[
\frac{P(X=k,Y=n-k)}{P(X+Y=n)}
\]

Substitute Poisson PMFs and simplify.

Final result:

\[
\binom{n}{k}\left(\frac{\lambda_1}{\lambda_1+\lambda_2}\right)^k
\left(\frac{\lambda_2}{\lambda_1+\lambda_2}\right)^{n-k}
\]

Thus conditional distribution is Binomial.

\section*{Q6. Transformation of Random Variable}

Given:

\[
X \sim \text{Uniform}[-2,2]
\]

Piecewise function:

\[
g(x)=x \quad [-2,-1]
\]

\[
g(x)=0 \quad (-1,1)
\]

\[
g(x)=x \quad [1,2]
\]

CDF Construction

\[
y<-2: 0
\]

\[
-2\le y<-1: \frac{y+2}{4}
\]

\[
-1\le y<0: \frac{1}{4}
\]

Jump at 0 of size $\frac{1}{2}$

\[
1\le y\le 2: \frac{y+2}{4}
\]

\[
y>2:1
\]

PDF:

For $(-2,-1)$ and $(1,2)$:

\[
f_Y(y)=\frac{1}{4}
\]

Point mass at 0:

\[
P(Y=0)=\frac{1}{2}
\]

\section*{Q7. Binomial Approximation (Stock Model)}

Price after 1000 periods:

\[
S_{1000}=S_0u^k d^{1000-k}
\]

Condition:

\[
u^k d^{1000-k}\ge1.30
\]

Taking logs:

\[
k(\ln u-\ln d)+1000\ln d\ge\ln(1.30)
\]

Solving gives:

\[
k\ge470
\]

Since:

\[
k\sim\text{Binomial}(1000,0.52)
\]

Required probability:

\[
\sum_{k=470}^{1000}\binom{1000}{k}(0.52)^k(0.48)^{1000-k}
\]

\section*{Q8. Exponential Distribution}

\[
\lambda=\frac{1}{2}
\]

PDF:

\[
\frac{1}{2}e^{-x/2}
\]

(a)

\[
P(X>2)=e^{-1}
\]

(b) Conditional Probability

\[
P(X\ge10|X>9)=\frac{e^{-5}}{e^{-9/2}}=e^{-1/2}
\]

\section*{Q9. Distribution of $Z$}

\[
Z=\frac{X^2+Y^2}{2}
\]

Start from CDF:

\[
F_Z(z)=P(X^2+Y^2\le2z)
\]

Since independent:

\[
f_{X,Y}(x,y)=f_X(x)f_Y(y)
\]

Region:

\[
x^2+y^2\le2z
\]

Double integral form derived.

Differentiate using Leibniz rule.

Final PDF:

\[
f_Z(z)=\int_{-z}^{z}f_X(x)\left[f_Y(\sqrt{2z-x^2})+f_Y(-\sqrt{2z-x^2})\right]dx
\]

\section*{Q10. Joint Distribution of $M_n$ and $M_{n+1}$}

\[
M_n=\max(X_1,\dots,X_n)
\]

Since i.i.d.:

\[
P(M_n\le x)=F(x)^n
\]

Two cases:

Case 1: $x\le y$

\[
P(M_n\le x,M_{n+1}\le y)=F(x)^nF(y)
\]

Case 2: $x>y$

\[
P(M_n\le x,M_{n+1}\le y)=F(y)^{n+1}
\]


\end{document}